%\documentclass[draft, 10pt]{beamer}
\documentclass[10pt]{beamer}
\usepackage[utf8]{inputenc}
\usepackage[T1]{fontenc}
\usepackage{lmodern}
\usepackage{graphicx}
\usepackage{booktabs}
\usepackage{bm}
\usepackage{tikz}
\usepackage{ulem}
\usepackage{appendixnumberbeamer}
\usepackage{hyperref}
\usepackage{ragged2e}
\setbeamersize{text margin left=0.5cm, text margin right=0.5cm}

\definecolor{unigreen}{RGB}{3,138,94}
\usecolortheme[named=unigreen]{structure}
\author[Braun and \"{O}zt\"{u}rk]{Sebastian T. Braun\textsuperscript{$\ast$}, Timur \"{O}zt\"{u}rk\textsuperscript{$\ddagger$}}
\institute[]{\textsuperscript{$\ast$}Kiel Institute for the World Economy, Kiel University, IZA \\ \vspace{1mm}
             \textsuperscript{$\ddagger$}University of Bayreuth}
\beamertemplatenavigationsymbolsempty

\graphicspath{
  {../GSWG2025/images/}
  {../../paper/figures/}
  {../../repository/regional/results/figures/}
}

%%%%%%%%%%%%%%%%%%%%%%%%%%%%%%%%%%%%%%%%%%%%%%%%%%%%%%%%%%%%%%%%%%%%
\begin{document}
\title[Confessional Politics]{From Weimar to West Germany:\\Confessional Politics, Party Fragmentation, and the Extreme Right}
\date[2026]{\today}
\setbeamercovered{dynamic}

% ============================================================
% TITLE
% ============================================================
\begin{frame}[plain,noframenumbering]
	\titlepage
\end{frame}

% ============================================================
\section{Introduction}
% ============================================================

\begin{frame}
	\frametitle{Motivation}
	\begin{itemize}
		\item Why did the Weimar Republic collapse, while the political system of the Federal Republic proved stable?
		\item Key structural difference: \textcolor{unigreen}{\textbf{fragmentation of the center-right}} in Weimar
		      \medskip
		\item \textbf{This paper:} Catholic-Protestant divisions drove that fragmentation---and the Nazi vote
		\item The CDU's postwar formation as an interdenominational \textit{Volkspartei} broke these dynamics
	\end{itemize}
\end{frame}

\begin{frame}
	\frametitle{This Paper}
	\begin{enumerate}
		\item \textbf{Quasi-experimental design} \\
		      {\small Compare adjacent Catholic and Protestant parishes at Reformation-era confessional boundaries}
		      \medskip
		\item \textbf{Novel parish-level data} \\
		      {\small 3,180 localities in Baden-Württemberg, 1919--1969 (vs.\ 104 county obs.\ in Falter data)}
		      \medskip
		\item \textbf{Survey mechanism evidence} \\
		      {\small 1968 Landtag survey: Catholics and Protestants had similar preferences; gap reflects church signals and church attachment}
	\end{enumerate}
\end{frame}

% ============================================================
\section{Historical Background}
% ============================================================

\begin{frame}
	\frametitle{Historical Background}
	\framesubtitle{Confessions in Weimar Germany, 1925}
	\begin{center}
		\includegraphics[width=0.72\linewidth]{ReligionWeimar1925.png}
	\end{center}
\end{frame}

\begin{frame}
	\frametitle{Historical Background}
	\framesubtitle{Reichstag elections, 1919--1933 \hfill \hyperlink{governments}{\beamergotobutton{Governments}}}
	\begin{center}
		\includegraphics[width=0.9\linewidth]{elections-weimar-1919-1933.pdf}
	\end{center}
\end{frame}

\begin{frame}
	\frametitle{Historical Background}
	\framesubtitle{Two very different center-right structures}
	\begin{columns}[T]
		\begin{column}{0.48\textwidth}
			\textbf{\textcolor{unigreen}{Catholics}}
			\begin{itemize}
				\item Strong \textit{Zentrum} voting norm
				\item Clergy warned against NSDAP
				\item Cohesive, well-anchored
			\end{itemize}
		\end{column}
		\begin{column}{0.48\textwidth}
			\textbf{\textcolor{unigreen}{Protestants}}
			\begin{itemize}
				\item No comparable organizing party
				\item Vote dispersed: DDP, DNVP, DVP, \ldots
				\item Structurally \textit{vulnerable} to Nazi consolidation
			\end{itemize}
		\end{column}
	\end{columns}
\end{frame}

% ============================================================
\section{Conceptual Framework}
% ============================================================

\begin{frame}
	\frametitle{Conceptual Framework}
	\framesubtitle{Spenkuch and Tillmann (2018) adapted to our setting}
	$$
		U_i^r(X)= \left\{\begin{array}{lr}
			-(X-t_i)^2+\lambda^r & \text{for } X = Z \text{ (Centre Party)}, \\
			-(X-t_i)^2           & \text{otherwise}
		\end{array}
		\right.
	$$
	\begin{itemize}
		\item Same preference distribution; differ only in norm: $\lambda^c > 0$, $\lambda^p < 0$
		\item \textbf{Testable:} any confessional gap reflects \textit{norms}, not preferences
	\end{itemize}
	\medskip
	\textbf{Hypotheses (Weimar):}
	\begin{enumerate}
		\item Fragmentation hump-shaped in Catholic share (1928); collapses in Protestant parishes by 1932
		\item Higher Catholic share $\Rightarrow$ lower NSDAP vote in 1932
	\end{enumerate}
	\textbf{Hypotheses (FRG):}
	\begin{enumerate}
		\item Fragmentation hump weakens after 1949 and dissolves by 1965 (CDU integrates denominations)
		\item Catholic parishes show lower NPD support in late 1960s, but gap much smaller than in 1932
	\end{enumerate}
\end{frame}

\begin{frame}
	\frametitle{Conceptual Framework: Weimar}
	\framesubtitle{Voting by voter type and confession}
	\begin{center}
		\includegraphics[width=1.0\linewidth]{voting-model-W.pdf}
	\end{center}
\end{frame}

\begin{frame}
	\frametitle{Conceptual Framework: Weimar}
	\framesubtitle{Predicted fragmentation by local Catholic share}
	\begin{center}
		\includegraphics[width=0.72\linewidth]{fragmentation-index-by-scenario.pdf}
	\end{center}
	Effective number of parties: $N = 1/\sum_j p_j^2$
\end{frame}

% ============================================================
\section{Empirical Strategy}
% ============================================================

\begin{frame}[label=identification]
	\frametitle{Empirical Strategy}
	\framesubtitle{\textit{Cuius regio, eius religio} --- comparing adjacent parishes}
	\begin{itemize}
		\item \textbf{Data:} 3,180 parishes (Baden and Württemberg), 1928--1969; 1968 Landtag survey ($N=895$)
		\item Post-1555 territorial fragmentation produced adjacent Catholic and Protestant parishes
	\end{itemize}
	\begin{center}
		\includegraphics[width=0.80\linewidth]{election-religion-map.png}
	\end{center}
\end{frame}

\begin{frame}
	\frametitle{Empirical Strategy}
	\framesubtitle{Identification strategies}
	\[y_{it} = \alpha + \beta\, CathShare_{i} + f(lat_i, lon_i) + \bm{X}_{it}\bm{\gamma} + \varepsilon_{it}\]
	\begin{enumerate}
		\item \textbf{Within-county:} county fixed effects (103 counties)
		\item \textbf{Within-neighborhood:} demean within 10 km radius (Mueller 2024; Becker 2025)
		\item \textbf{Religious enclaves:} 44 enclave cases vs.\ surrounding parishes \hyperlink{enclaves}{\beamergotobutton{Example}}
	\end{enumerate}
	\medskip
	Large raw differences between Catholic and Protestant parishes vanish within counties and neighborhoods. \hyperlink{balancing}{\beamergotobutton{Balancing table}}
\end{frame}

% ============================================================
\section{Results: Weimar Republic}
% ============================================================

\begin{frame}
	\frametitle{Fragmentation in Weimar}
	\framesubtitle{Baden, May 1928 and November 1932}
	\begin{center}
		\includegraphics[width=1.0\linewidth]{fragmentation-baden-1928-1932.pdf}
	\end{center}
	\begin{itemize}
		\item \textbf{1928:} Inverted U-shape confirmed
		\item \textbf{1932:} NSDAP consolidated Protestant vote; fragmentation collapsed in Protestant/mixed parishes
	\end{itemize}
\end{frame}

\begin{frame}[label=resultsNSDAP]
	\frametitle{Nazi Vote, November 1932}
	\framesubtitle{Baden --- OLS regressions}
	\makebox[1.0\linewidth][c]{
		\begin{footnotesize}
			\begin{tabular}{lcccc}
				\toprule
				               & (1)         & (2)         & (3)         & (4)         \\
				               & Baseline    & County FE   & Radius      & Enclaves    \\
				\midrule
				Catholics (\%) & $-$0.465*** & $-$0.496*** & $-$0.516*** & $-$0.486*** \\
				               & (0.027)     & (0.030)     & (0.028)     & (0.072)     \\
				\midrule
				Observations   & 1,436       & 1,380       & 1,380       & 385         \\
				$R^2$          & 0.55        & 0.66        & 0.56        & 0.72        \\
				Controls       & no          & yes         & yes         & yes         \\
				\bottomrule
				\multicolumn{5}{l}{\scriptsize County-clustered (enclave-clustered in col.\ 4) SE. ***, ** = 1\%, 5\%.}
			\end{tabular}
		\end{footnotesize}
	}
	\begin{itemize}
		\item All-Catholic vs.\ all-Protestant: $\approx$\textbf{47 pp} lower Nazi vote
		\item Catholic share alone explains \textbf{53\%} of cross-parish variation \hyperlink{fullNSDAP}{\beamergotobutton{Full table}}
	\end{itemize}
\end{frame}

\begin{frame}
	\frametitle{How Important is Religion?}
	\framesubtitle{Dominance decomposition of $R^2$ --- Baden, November 1932}
	\begin{center}
		\includegraphics[width=0.72\linewidth]{dominance.pdf}
	\end{center}
	Catholic share accounts for $\approx$\textbf{75\%} of the model's explained variation.
\end{frame}

% ============================================================
\section{Results: Federal Republic}
% ============================================================

\begin{frame}
	\frametitle{Fragmentation in the Federal Republic}
	\framesubtitle{Baden-Württemberg, August 1949 and September 1965}
	\begin{center}
		\includegraphics[width=1.0\linewidth]{fragmentation-bw-panel-cath1950.pdf}
	\end{center}
	\begin{itemize}
		\item \textbf{1949:} Hump persists but weaker; \textbf{1965:} largely dissolved
		\item Largest reductions in mixed and Protestant parishes --- CDU's integrative role
	\end{itemize}
\end{frame}

\begin{frame}
	\frametitle{Centre Party (1928) vs.\ CDU (1965)}
	\framesubtitle{Baden --- CDU gained most ground in Protestant parishes}
	\begin{center}
		\includegraphics[width=0.9\linewidth]{vote-share-zentrum-cdu-change-panel-baden.pdf}
	\end{center}
\end{frame}

\begin{frame}
	\frametitle{NPD Support in the Late 1960s}
	\framesubtitle{Protestant areas remained more susceptible}
	\makebox[1.0\linewidth][c]{
		\begin{footnotesize}
			\begin{tabular}{lcccc}
				\toprule
				               & \multicolumn{2}{c}{CDU} & \multicolumn{2}{c}{NPD}                              \\
				               & 1949                    & 1969                    & 1968 state  & 1969 federal \\
				\midrule
				Catholics (\%) & 0.497***                & 0.378***                & $-$0.084*** & $-$0.059***  \\
				               & (0.023)                 & (0.019)                 & (0.011)     & (0.012)      \\
				\midrule
				Obs.           & 1,261                   & 3,117                   & 1,331       & 3,117        \\
				Within transf. & radius                  & radius                  & radius      & radius       \\
				\bottomrule
				\multicolumn{5}{l}{\scriptsize Controls included. County-clustered SE.}
			\end{tabular}
		\end{footnotesize}
	}
	\begin{itemize}
		\item CDU stronger in Catholic parishes throughout, but gap narrowed
		\item NPD gap much smaller than 1932 Nazi vote gap ($-$0.06 vs.\ $-$0.47)
	\end{itemize}
\end{frame}

% ============================================================
\section{Individual-Level Survey Evidence}
% ============================================================

\begin{frame}
	\frametitle{NPD Support by Confession}
	\framesubtitle{1968 Baden-Württemberg Landtag survey (ZA0327, $N=895$)}
	\begin{center}
		\includegraphics[width=0.82\linewidth]{NPD_vote_robindicators_by_confession.pdf}
	\end{center}
	Protestants substantially more likely to support NPD across all four indicators.
\end{frame}

\begin{frame}
	\frametitle{Similar Underlying Political Preferences}
	\begin{center}
		\includegraphics[width=0.82\linewidth]{agreement_with_NPD_positions_by_confession.pdf}
	\end{center}
	Agreement with NPD positions: \textbf{virtually identical} across confessions $\rightarrow$ preference explanation ruled out.
\end{frame}

\begin{frame}
	\frametitle{Church Signals on CDU and NPD}
	\begin{center}
		\includegraphics[width=0.82\linewidth]{perception_CDU_NPD_by_church.pdf}
	\end{center}
	\begin{itemize}
		\item 86\% of Catholics: Catholic Church's stance on NPD is ``bad/very bad''
		\item 24\% of Protestants: Protestant Church's stance on NPD is neutral or positive
	\end{itemize}
\end{frame}

\begin{frame}
	\frametitle{Church Importance and Voting}
	\begin{center}
		\includegraphics[width=0.82\linewidth]{LTW1968_CDU_NPD_vote_by_church_importance.pdf}
	\end{center}
	\begin{itemize}
		\item Among the highly devout: NPD support \textit{comparably low} in both confessions
		\item But 29\% of Catholics vs.\ 18\% of Protestants are in that group
	\end{itemize}
\end{frame}

% ============================================================
\section{Conclusion}
% ============================================================

\begin{frame}
	\frametitle{Conclusion}
	\begin{itemize}
		\item \textbf{Weimar:} Confessional divisions fragmented the center-right; Catholic share alone explains $>$50\% of the 1932 Nazi vote
		      \medskip
		\item \textbf{Federal Republic:} CDU reduced denominational fragmentation substantially; Protestant areas remained more exposed to the NPD
		      \medskip
		\item \textbf{Mechanism:} Confessional gap reflects elite influence (church signals + church attachment), not differences in preferences
		      \medskip
		\item \textbf{General:} Cross-cleavage parties can mitigate fragmentation risks generated by deep social divisions
	\end{itemize}
	\vspace{0.8em}
	\begin{center}
		{\large Thank you!} \quad
		sebastian.braun@kielinstitut.de \quad $\cdot$ \quad timur.oeztuerk@uni-bayreuth.de
	\end{center}
\end{frame}

% ============================================================
\appendix
% ============================================================

\begin{frame}[label=governments, noframenumbering]
	\frametitle{Weimar Governments, 1919--1933}
	\framesubtitle{Source: Lehmann (2010)}
	\begin{center}
		\includegraphics[width=0.65\linewidth]{governments.pdf}
	\end{center}
\end{frame}

\begin{frame}[label=enclaves, noframenumbering]
	\frametitle{Religious Enclaves: Example}
	\framesubtitle{Emmendingen district, 1932}
	\begin{center}
		\includegraphics[width=1.0\linewidth]{election-religion-map-enclave_OAEmmendingen.png}
	\end{center}
	\begin{flushright}
		\hyperlink{identification}{\beamergotobutton{Back}}
	\end{flushright}
\end{frame}

\begin{frame}[label=balancing, noframenumbering]
	\frametitle{Balancing Tests}
	\makebox[\linewidth][c]{
		\begin{scriptsize}
			\begin{tabular}{lccc}
				\toprule
				                       & Unconditional & Within county & Within 10 km \\
				\midrule
				Average elevation (m)  & 104.9***      & 10.2          & $-$0.8       \\
				Broad river valleys    & 0.071***      & 0.006         & 0.004        \\
				Loess area             & $-$0.167***   & $-$0.068**    & $-$0.040**   \\
				Population 1925 (logs) & $-$0.190**    & $-$0.047      & $-$0.063     \\
				Industry share (\%)    & $-$10.7***    & 1.16          & 0.31         \\
				\bottomrule
			\end{tabular}
		\end{scriptsize}
	}
	\begin{flushright}
		\hyperlink{identification}{\beamergotobutton{Back}}
	\end{flushright}
\end{frame}

\begin{frame}[label=fullNSDAP, noframenumbering]
	\frametitle{Nazi Vote: Full Regression Table}
	\framesubtitle{Baden, November 1932}
	\makebox[1.0\linewidth][c]{
		\begin{tiny}
			\begin{tabular}{lccccccc}
				\toprule
				               & (1)         & (2)         & (3)         & (4)         & (5)         & (6)         & (7)             \\
				               & Baseline    & + Spatial   & + Controls  & County FE   & Radius      & Enclaves    & $\Delta$ 28--32 \\
				\midrule
				Catholics (\%) & $-$0.465*** & $-$0.454*** & $-$0.496*** & $-$0.516*** & $-$0.486*** & $-$0.534*** & $-$0.432***     \\
				               & (0.027)     & (0.026)     & (0.030)     & (0.028)     & (0.025)     & (0.072)     & (0.026)         \\
				\midrule
				Obs.           & 1,436       & 1,380       & 1,380       & 1,380       & 1,380       & 385         & 1,380           \\
				$R^2$          & 0.55        & 0.57        & 0.62        & 0.66        & 0.56        & 0.72        & 0.46            \\
				Controls       & no          & no          & yes         & yes         & yes         & yes         & yes             \\
				\bottomrule
			\end{tabular}
		\end{tiny}
	}
	\begin{flushright}
		\hyperlink{resultsNSDAP}{\beamergotobutton{Back}}
	\end{flushright}
\end{frame}

\begin{frame}[noframenumbering]
	\frametitle{Individual-Level Regressions}
	\framesubtitle{Determinants of CDU and NPD voting, 1968}
	\makebox[1.0\linewidth][c]{
		\begin{tiny}
			\begin{tabular}{lcccccc}
				\toprule
				                                  & \multicolumn{3}{c}{CDU vote} & \multicolumn{3}{c}{NPD vote}                                                     \\
				                                  & (1)                          & (2)                          & (3)      & (4)        & (5)         & (6)         \\
				\midrule
				Catholic                          & 0.197***                     & 0.177***                     & 0.131**  & $-$0.038** & $-$0.030*   & $-$0.027    \\
				Church very important             &                              & 0.211***                     & 0.182*** &            & $-$0.044*** & $-$0.041*** \\
				Catholic $\times$ Church v.\ imp. &                              &                              & 0.079    &            &             & $-$0.005    \\
				\midrule
				Obs.                              & 587                          & 587                          & 587      & 587        & 587         & 587         \\
				\bottomrule
				\multicolumn{7}{l}{\scriptsize Controls: age, gender, education, marital status, community size.}
			\end{tabular}
		\end{tiny}
	}
\end{frame}

\end{document}
